% ================================================================
%  KR260 / DAPHNE-V3 – DAPHNE-15 Golden Recovery (2026-03-12)
% ================================================================

\documentclass[a4paper,11pt]{article}

% ---------- packages ------------------------------------------------
\usepackage[utf8]{inputenc}
\usepackage[T1]{fontenc}
\usepackage{lmodern}
\usepackage{geometry}
\usepackage{hyperref}
\usepackage{enumitem}
\usepackage{xcolor}
\usepackage{graphicx}
\usepackage{booktabs}
\usepackage{caption}
\usepackage{float}
\usepackage{minted}
\usepackage{amsmath}

\geometry{margin=25mm}

% ---------- minted defaults -----------------------------------------
\setminted{
  fontsize   = \small,
  frame      = single,
  breaklines = true,
  style      = colorful
}

\title{KR260 / DAPHNE-V3\\[4pt]\large DAPHNE-15 Golden Recovery (2026-03-12)}
\author{}
\date{\today}

\begin{document}
\maketitle
\tableofcontents
\clearpage

%--------------------------------------------------------------------
\section{Overview}

This runbook documents the recovery path that worked on March 12, 2026 for
\texttt{daphne-15}.

The board boot chain is:
\[
\textbf{QSPI (BOOT.BIN)} \rightarrow \textbf{FSBL/PMUFW/ATF} \rightarrow \textbf{U-Boot} \rightarrow \textbf{Linux on eMMC}.
\]

The recovery goal is:
\begin{enumerate}[label=\arabic*.]
  \item reach a working \texttt{ZynqMP>} prompt from QSPI or JTAG,
  \item replace the broken eMMC DTB with the known-good golden DTB,
  \item boot a \textbf{true RAM shell} (\texttt{/ = rootfs}, not \texttt{/dev/mmcblk0p2}),
  \item bring up Ethernet with the \texttt{daphne-15} identity,
  \item overwrite \texttt{mmcblk0p2} with the golden rootfs image,
  \item rewrite \texttt{mmcblk0p1} boot files from the golden bundle,
  \item reboot normally and re-pin \texttt{daphne-15} identity/services in Linux.
\end{enumerate}

\paragraph{Golden source used in this procedure.}
\begin{itemize}
  \item Golden bundle: \texttt{/nfs/home/marroyav/golden/daphne14-2026-03-12}
  \item Working JTAG boot script: \texttt{/nfs/home/marroyav/daphne_v3/linux_13/boot.tcl}
\end{itemize}

\paragraph{Important conclusion from the recovery.}
The board could reach U-Boot from QSPI, so QSPI was \textbf{not} reflashed first.
The actual blocker was the installed DTB on \texttt{mmc0:1}. Once the golden DTB
was written to \texttt{mmc0:1}, the usable management uplink appeared on
\texttt{eth0}.

\clearpage
%--------------------------------------------------------------------
\section{Artifacts and Known-Good Files}

\subsection*{Golden bundle}
\begin{listing}[H]
\begin{minted}{bash}
/nfs/home/marroyav/golden/daphne14-2026-03-12/
├── boot/
│   ├── BOOT.BIN
│   ├── Image
│   ├── boot.scr
│   ├── ramdisk.cpio.gz.u-boot
│   ├── system.dtb
│   └── system-zynqmp-sck-kr-g-revB.dtb
├── mmcblk0p2.ext4.img.gz
├── qspi/
└── meta/
\end{minted}
\caption{Golden recovery bundle}
\end{listing}

\subsection*{Golden DTB size}
The correct DTB size used with \texttt{fatwrite} was:
\begin{listing}[H]
\begin{minted}{bash}
wc -c /nfs/home/marroyav/golden/daphne14-2026-03-12/boot/system.dtb
# 49280 .../system.dtb
\end{minted}
\caption{Golden DTB size}
\end{listing}

In U-Boot, this size must be written as hexadecimal:
\begin{listing}[H]
\begin{minted}{text}
49280 decimal = 0xc080 hex
\end{minted}
\caption{DTB size for U-Boot fatwrite}
\end{listing}

\clearpage
%--------------------------------------------------------------------
\section{Consoles}

\subsection*{Serial console}
\begin{listing}[H]
\begin{minted}{bash}
screen -D -RR -S daphne-15 /dev/ttyUSB2 115200
\end{minted}
\caption{Serial console to daphne-15}
\end{listing}

\subsection*{XSCT / JTAG console}
The reliable JTAG bootstrap path used the existing script in
\texttt{linux\_13}:
\begin{listing}[H]
\begin{minted}{bash}
cd /nfs/home/marroyav/daphne_v3/linux_13
source /nfs/sw/fpga/Xilinx/Vitis/2024.1/settings64.sh
xsct boot.tcl
\end{minted}
\caption{Use the known-good boot.tcl to reach U-Boot via JTAG}
\end{listing}

\paragraph{When to use XSCT.}
Use \texttt{boot.tcl} only if the board does not reach a usable
\texttt{ZynqMP>} prompt from QSPI on its own, or if a fresh JTAG-controlled
start is desired before patching eMMC.

\clearpage
%--------------------------------------------------------------------
\section{Patch the Installed DTB on \texttt{mmc0:1}}

\subsection*{Why this step comes first}
The failed recovery attempts showed that:
\begin{itemize}
  \item the board could boot a RAM shell,
  \item but Ethernet did not work correctly with the installed DTB,
  \item after replacing \texttt{system.dtb} on \texttt{mmc0:1} with the golden
        DTB, the usable uplink appeared on \texttt{eth0}.
\end{itemize}

\subsection*{Stage the golden DTB into DDR from XSCT}
On \texttt{np04-onl-004}:
\begin{listing}[H]
\begin{minted}{bash}
cd /nfs/home/marroyav/daphne_v3/linux_13
xsct
\end{minted}
\caption{Start XSCT}
\end{listing}

Inside the \texttt{xsct} shell:
\begin{listing}[H]
\begin{minted}{tcl}
source boot.tcl
targets -set -nocase -filter {name =~ "*A53*#0"}
stop
dow -data /nfs/home/marroyav/golden/daphne14-2026-03-12/boot/system.dtb 0x00100000
con
exit
\end{minted}
\caption{Stage the golden DTB at 0x00100000}
\end{listing}

\subsection*{Write the staged DTB into eMMC p1 from U-Boot}
At the serial \texttt{ZynqMP>} prompt:
\begin{listing}[H]
\begin{minted}{bash}
fdt addr 0x00100000
fdt header
mmc dev 0
fatwrite mmc 0:1 0x00100000 system.dtb 0xc080
fatwrite mmc 0:1 0x00100000 system-zynqmp-sck-kr-g-revB.dtb 0xc080
fatls mmc 0:1
\end{minted}
\caption{Rewrite the installed DTB with the golden one}
\end{listing}

\paragraph{Expected result.}
\texttt{fatls mmc 0:1} should show:
\begin{listing}[H]
\begin{minted}{text}
49280   system.dtb
49280   system-zynqmp-sck-kr-g-revB.dtb
\end{minted}
\caption{Correct DTB file sizes on eMMC p1}
\end{listing}

\clearpage
%--------------------------------------------------------------------
\section{Boot a True RAM Shell}

\subsection*{Goal}
Boot Linux from:
\begin{itemize}
  \item \texttt{Image} from \texttt{mmc0:1},
  \item \texttt{ramdisk.cpio.gz.u-boot} from \texttt{mmc0:1},
  \item the newly written golden \texttt{system.dtb},
\end{itemize}
with \texttt{rdinit=/bin/sh} so that the root filesystem remains the
initramfs (\texttt{rootfs / rootfs}) and not \texttt{/dev/mmcblk0p2}.

\subsection*{U-Boot commands}
\begin{listing}[H]
\begin{minted}{bash}
mmc dev 0
fatload mmc 0:1 ${kernel_addr_r} Image
fatload mmc 0:1 ${fdt_addr_r} system.dtb
fatload mmc 0:1 ${ramdisk_addr_r} ramdisk.cpio.gz.u-boot
setenv bootargs 'console=ttyPS1,115200 earlycon root=/dev/ram rw rdinit=/bin/sh devtmpfs.mount=1'
booti ${kernel_addr_r} ${ramdisk_addr_r} ${fdt_addr_r}
\end{minted}
\caption{Boot to a real RAM shell}
\end{listing}

\subsection*{Verify the shell is running from RAM}
At the shell prompt (\texttt{/ \#}):
\begin{listing}[H]
\begin{minted}{bash}
mkdir -p /proc /sys /dev
mount -t proc proc /proc
mount -t sysfs sysfs /sys
mount -t devtmpfs devtmpfs /dev 2>/dev/null || true
cat /proc/mounts
\end{minted}
\caption{Verify the root is the initramfs}
\end{listing}

\paragraph{Expected result.}
The output must include:
\begin{listing}[H]
\begin{minted}{text}
rootfs / rootfs rw 0 0
\end{minted}
\caption{Correct RAM-root signature}
\end{listing}

and must \textbf{not} show:
\begin{listing}[H]
\begin{minted}{text}
/dev/mmcblk0p2 on /
\end{minted}
\caption{Incorrect result: not safe for flashing}
\end{listing}

\clearpage
%--------------------------------------------------------------------
\section{Bring Up the Network in the RAM Shell}

\subsection*{What changed after the DTB fix}
With the corrected golden DTB:
\begin{itemize}
  \item \texttt{eth0} became the usable management uplink,
  \item \texttt{eth0} corresponded to \texttt{ff0b0000} and came up with
        \texttt{fixed/sgmii},
  \item the working \texttt{daphne-15} recovery address was applied on
        \texttt{eth0}.
\end{itemize}

\subsection*{Bring the interfaces up}
\begin{listing}[H]
\begin{minted}{bash}
ip link set eth0 up 2>/dev/null || true
ip link set eth1 up 2>/dev/null || true
\end{minted}
\caption{Bring the links up}
\end{listing}

\subsection*{Configure the recovery uplink on eth0}
\begin{listing}[H]
\begin{minted}{bash}
ip addr flush dev eth0 2>/dev/null || true
ip addr flush dev eth1 2>/dev/null || true
ip route del default 2>/dev/null || true

ip link set eth0 down
ip link set eth0 address ba:be:ba:d1:cc:ff
ip link set eth0 up

ip addr add 10.73.137.16/24 dev eth0
ip route replace default via 10.73.137.1 dev eth0

ip -4 addr show dev eth0
ip route
ping -c 3 10.73.137.1
ping -c 3 10.73.138.64
\end{minted}
\caption{Configure the working daphne-15 recovery uplink}
\end{listing}

\paragraph{Expected result.}
Both pings succeeded after the DTB replacement:
\begin{itemize}
  \item gateway: \texttt{10.73.137.1}
  \item bridge host: \texttt{10.73.138.64}
\end{itemize}

\clearpage
%--------------------------------------------------------------------
\section{Overwrite \texttt{mmcblk0p2} from the Golden Rootfs}

\subsection*{Sender on \texttt{np04-onl-004}}
\begin{listing}[H]
\begin{minted}{bash}
GOLDEN=/nfs/home/marroyav/golden/daphne14-2026-03-12
gzip -dc $GOLDEN/mmcblk0p2.ext4.img.gz > $GOLDEN/mmcblk0p2.ext4
dd if=$GOLDEN/mmcblk0p2.ext4 bs=1M status=progress | nc -l 9002
\end{minted}
\caption{Serve the golden rootfs image}
\end{listing}

\subsection*{Receiver on \texttt{daphne-15} RAM shell}
\begin{listing}[H]
\begin{minted}{bash}
nc 10.73.138.64 9002 | dd of=/dev/mmcblk0p2 bs=1M
sync
\end{minted}
\caption{Overwrite the eMMC rootfs partition}
\end{listing}

\paragraph{Why this is safe.}
Because the board is running from the initramfs (\texttt{rootfs / rootfs}) and
not from \texttt{/dev/mmcblk0p2}, the rootfs partition can be replaced in-place.

\paragraph{BusyBox note.}
The RAM shell used during this recovery provides the BusyBox implementation of
\texttt{dd}. It does \textbf{not} support \texttt{conv=fsync} or
\texttt{status=progress}. The transfer is considered active while:
\begin{itemize}
  \item the \texttt{dd} command on \texttt{daphne-15} has not returned to the prompt,
  \item the sender command on \texttt{np04-onl-004}
        (\texttt{dd if=... status=progress | nc -l 9002})
        is still running.
\end{itemize}
\texttt{nc} itself does not provide progress output. Progress is shown on the
sender by GNU \texttt{dd}.

\clearpage
%--------------------------------------------------------------------
\section{Rewrite \texttt{mmcblk0p1} Boot Files from the Golden Bundle}

\subsection*{Mount the FAT partition on \texttt{daphne-15}}
\begin{listing}[H]
\begin{minted}{bash}
mkdir -p /mnt/p1
mount -t vfat /dev/mmcblk0p1 /mnt/p1
\end{minted}
\caption{Mount eMMC boot partition}
\end{listing}

\subsection*{Send files one-by-one from \texttt{np04-onl-004}}
\begin{listing}[H]
\begin{minted}{bash}
GOLDEN=/nfs/home/marroyav/golden/daphne14-2026-03-12
nc -l 9003 < $GOLDEN/boot/BOOT.BIN
nc -l 9004 < $GOLDEN/boot/Image
nc -l 9005 < $GOLDEN/boot/boot.scr
nc -l 9006 < $GOLDEN/boot/ramdisk.cpio.gz.u-boot
\end{minted}
\caption{Serve the golden boot files}
\end{listing}

\subsection*{Receive files on \texttt{daphne-15}}
\begin{listing}[H]
\begin{minted}{bash}
nc 10.73.138.64 9003 > /mnt/p1/BOOT.BIN
nc 10.73.138.64 9004 > /mnt/p1/Image
nc 10.73.138.64 9005 > /mnt/p1/boot.scr
nc 10.73.138.64 9006 > /mnt/p1/ramdisk.cpio.gz.u-boot
sync
umount /mnt/p1
\end{minted}
\caption{Update the eMMC boot partition}
\end{listing}

\paragraph{Note.}
The DTB was already rewritten earlier from U-Boot with:
\begin{itemize}
  \item \texttt{system.dtb}
  \item \texttt{system-zynqmp-sck-kr-g-revB.dtb}
\end{itemize}
so only the remaining boot files need to be transferred here.

\clearpage
%--------------------------------------------------------------------
\section{Reboot and Re-pin DAPHNE-15 Identity}

\subsection*{Reboot}
\begin{listing}[H]
\begin{minted}{bash}
reboot -f
\end{minted}
\caption{Reboot after eMMC rewrite}
\end{listing}

\paragraph{Observed result during the 2026-03-12 recovery.}
After the successful golden eMMC rewrite, the board booted fully into Linux but
identified itself as the cloned source board:
\begin{listing}[H]
\begin{minted}{text}
PetaLinux 2024.1+release-S05201002 NP04-DAPHNE-014.CERN.CH ttyPS1
\end{minted}
\caption{Expected first boot before re-pinning daphne-15 identity}
\end{listing}

\subsection*{Apply the board-specific persistent identity}
After the board boots normally into Linux again:
\begin{listing}[H]
\begin{minted}{bash}
cd ~/daphne-server/utils
sudo FW_APP=MEZ_SELF_TRIG_V15_OL_UPGRADED ./bootstrap_blank_board.sh daphne-15
sudo reboot
\end{minted}
\caption{Install the persistent daphne-15 identity}
\end{listing}

\paragraph{What this fixes persistently.}
\begin{itemize}
  \item changes the cloned \texttt{daphne-14} identity into \texttt{daphne-15},
  \item hostname,
  \item deterministic ff0b/ff0c MAC and IP configuration,
  \item DNS and NTP,
  \item firmware service,
  \item Hermes / DAPHNE stack startup path.
\end{itemize}

\clearpage
%--------------------------------------------------------------------
\section{Verification}

\subsection*{From U-Boot}
\begin{listing}[H]
\begin{minted}{bash}
mmc dev 0
fatls mmc 0:1
printenv bootcmd bootdelay
\end{minted}
\caption{Basic U-Boot verification}
\end{listing}

\subsection*{From Linux after final boot}
\begin{listing}[H]
\begin{minted}{bash}
hostname
ip -4 addr
ip route
systemctl status force-ff0b-net.service --no-pager
systemctl status firmware.service --no-pager
xmutil listapps
\end{minted}
\caption{Final Linux verification}
\end{listing}

\paragraph{Expected identity before and after bootstrap.}
\begin{itemize}
  \item Before running \texttt{bootstrap\_blank\_board.sh daphne-15}: the system may
        still present itself as \texttt{NP04-DAPHNE-014.CERN.CH}.
  \item After running \texttt{bootstrap\_blank\_board.sh daphne-15} and rebooting:
        the hostname/network identity should match \texttt{daphne-15}.
\end{itemize}

\clearpage
%--------------------------------------------------------------------
\section{Troubleshooting Notes from the 2026-03-12 Recovery}

\begin{description}
  \item[\texttt{Bad Linux ARM64 Image magic!}] The kernel image placed in DDR by JTAG was overwritten before \texttt{booti}. Do not rely on the JTAG-loaded kernel path for this recovery.
  \item[\texttt{root=/dev/ram} but \texttt{/} still becomes \texttt{/dev/mmcblk0p2}] The initramfs pivoted to eMMC. Use \texttt{rdinit=/bin/sh} and verify \texttt{rootfs / rootfs} before flashing.
  \item[\texttt{Network is unreachable}] If the DTB is wrong, the board may show link without a usable network path. Fix \texttt{system.dtb} on \texttt{mmc0:1} first.
  \item[\texttt{fatwrite} wrote the wrong size] U-Boot interpreted the numeric argument as hexadecimal. For the 49\,280-byte DTB, use \texttt{0xc080}, not \texttt{49280}.
  \item[\texttt{loady} over serial timing out] This path was not needed once the DTB was staged through XSCT and written from U-Boot.
\end{description}

\clearpage
%--------------------------------------------------------------------
\section*{Appendix: Quick Commands}

\begin{listing}[H]
\begin{minted}[linenos]{bash}
# --- Stage golden DTB in DDR from XSCT ---
cd /nfs/home/marroyav/daphne_v3/linux_13
xsct
source boot.tcl
targets -set -nocase -filter {name =~ "*A53*#0"}
stop
dow -data /nfs/home/marroyav/golden/daphne14-2026-03-12/boot/system.dtb 0x00100000
con
exit

# --- Write the golden DTB into mmc0:1 from U-Boot ---
fdt addr 0x00100000
fdt header
mmc dev 0
fatwrite mmc 0:1 0x00100000 system.dtb 0xc080
fatwrite mmc 0:1 0x00100000 system-zynqmp-sck-kr-g-revB.dtb 0xc080

# --- Boot the true RAM shell ---
fatload mmc 0:1 ${kernel_addr_r} Image
fatload mmc 0:1 ${fdt_addr_r} system.dtb
fatload mmc 0:1 ${ramdisk_addr_r} ramdisk.cpio.gz.u-boot
setenv bootargs 'console=ttyPS1,115200 earlycon root=/dev/ram rw rdinit=/bin/sh devtmpfs.mount=1'
booti ${kernel_addr_r} ${ramdisk_addr_r} ${fdt_addr_r}

# --- Configure the working uplink ---
ip link set eth0 down
ip link set eth0 address ba:be:ba:d1:cc:ff
ip link set eth0 up
ip addr add 10.73.137.16/24 dev eth0
ip route replace default via 10.73.137.1 dev eth0

# --- Overwrite rootfs ---
gzip -dc /nfs/home/marroyav/golden/daphne14-2026-03-12/mmcblk0p2.ext4.img.gz \
  > /nfs/home/marroyav/golden/daphne14-2026-03-12/mmcblk0p2.ext4
dd if=/nfs/home/marroyav/golden/daphne14-2026-03-12/mmcblk0p2.ext4 \
  bs=1M status=progress | nc -l 9002

# On daphne-15:
nc 10.73.138.64 9002 | dd of=/dev/mmcblk0p2 bs=1M

# --- Rewrite /boot from the golden bundle ---
mount -t vfat /dev/mmcblk0p1 /mnt/p1
nc 10.73.138.64 9003 > /mnt/p1/BOOT.BIN
nc 10.73.138.64 9004 > /mnt/p1/Image
nc 10.73.138.64 9005 > /mnt/p1/boot.scr
nc 10.73.138.64 9006 > /mnt/p1/ramdisk.cpio.gz.u-boot
sync
umount /mnt/p1
\end{minted}
\caption{Everything needed for the 2026-03-12 recovery}
\end{listing}

\end{document}
